\chapter{Study Sheets on Scientific Atheism}

\begingroup

%% ── EPIGRAPH ────────────────────────────────
\begin{quote}
\textit{``The whole problem with the world is that fools and fanatics are always
so certain of themselves, and wiser people so full of doubts.''}

\medskip
\hfill--- Bertrand Russell
\end{quote}

\begin{quote}
\textit{``Is God willing to prevent evil, but not able? Then he is not
omnipotent.  Is he able, but not willing?  Then he is malevolent.
Is he both able and willing?  Then whence cometh evil?
Is he neither able nor willing?  Then why call him God?''}

\medskip
\hfill--- Epicurus (attr.)
\end{quote}

%%  I — GENERAL THESES
\section{General Theses: The Epistemology of Disbelief}

\subsection{Sheet 1: The Burden of Proof}
The foundational principle of scientific atheism is not that ``God does not
exist'' as a dogmatic assertion, but that \textbf{extraordinary claims require
extraordinary evidence}.

\begin{enumerate}
  \item \textbf{Russell's Teapot.}
    If someone claims that a china teapot orbits the Sun between Earth and Mars,
    the burden of proof rests on the claimant, not on the sceptic.  Religious
    claims about a transcendent being are no different epistemologically.
  \item \textbf{The null hypothesis.}
    The rational default is non-belief in entities for which no falsifiable
    evidence is presented.
  \item \textbf{Unfalsifiability as a defect.}
    Theological propositions are typically constructed so that \emph{no
    conceivable observation} could refute them.  This makes them
    epistemically vacuous, not unassailable.
\end{enumerate}

\subsection{Sheet 2: Contradictions in the Concept of God}

Classical theism attributes three properties to God simultaneously:
\textbf{omnipotence}, \textbf{omniscience}, and \textbf{omnibenevolence}.
Each pair generates well-known paradoxes:

\begin{itemize}
  \item \textbf{The Omnipotence Paradox.}
    Can God create a stone so heavy that He cannot lift it?  Either answer
    limits omnipotence.
  \item \textbf{The Problem of Evil (Theodicy).}
    If God is all-knowing, all-powerful, and all-good, the existence of
    gratuitous suffering is logically inexplicable.  Classical ``free will''
    defences fail to account for natural evil
    (earthquakes, childhood leukaemia, pandemics).
  \item \textbf{Divine foreknowledge vs.\ free will.}
    If God knows all future events with certainty, then the future is fixed,
    and human freedom is illusory.  This undermines the entire apparatus of
    sin, merit, and punishment.
\end{itemize}

\subsection{Sheet 3: Psychological \& Sociological Origins of Religion}

\begin{itemize}
  \item \textbf{Feuerbach's projection thesis.}
    God is an externalisation of idealised human qualities.  Theology is
    anthropology in disguise.
  \item \textbf{Marx's social function.}
    Religion is the ``opium of the people''---it
    anaesthetises the exploited against the perception of their own condition,
    serving the interests of ruling classes.
  \item \textbf{Durkheim's social cohesion.}
    The ``sacred'' is a symbolic representation of society itself.  Ritual
    reinforces group solidarity, not metaphysical truth.
  \item \textbf{Cognitive science.}
    Humans exhibit an evolved \emph{hyperactive agency detection
    device} (HADD): we over-attribute intentionality to
    natural phenomena---a survival trait that predisposes us to theism.
\end{itemize}

\paragraph{Core Methodological Principle}
Scientific atheism does not assert certainty about metaphysical negatives.
It insists that belief should be proportioned to evidence, and that where
evidence is absent, suspension of judgement---not faith---is the rational
stance.

%%  II — MONOTHEISM
\section{Monotheism: Critique of the One-God Concept}

\subsection{Sheet 4: Historical Contingency of Monotheism}

Monotheism is not a ``natural'' or inevitable theological development.  It
emerged from specific historical pressures:

\begin{enumerate}
  \item \textbf{From polytheism to monolatry.}
    Early Israelite religion was not monotheistic but
    \emph{monolatrous}: Yahweh was
    worshipped as the \emph{national} god among many existing gods
    (cf.\ Psalm~82, Deuteronomy~32:8--9 in the Dead Sea Scrolls reading).
  \item \textbf{Political consolidation.}
    The Josiah reforms (c.~621~BCE) centralised cult
    worship in Jerusalem, suppressing rival sanctuaries and local deities.
    Monotheism was, in part, a political technology for national unification.
  \item \textbf{Persian influence.}
    Zoroastrian dualism profoundly shaped post-exilic
    Jewish theology: eschatology, angelology, the figure
    of Satan, and bodily resurrection.
\end{enumerate}

\subsection{Sheet 5: Logical Problems Specific to Monotheism}

\begin{itemize}
  \item \textbf{The problem of divine attributes.}
    A ``simple'' God (Aquinas) who is identical with His
    attributes cannot be a person in any meaningful sense---persons have
    contingent properties.  If God is love itself, He cannot \emph{choose}
    to love; love becomes a tautology.
  \item \textbf{Creation \emph{ex nihilo}.}
    If God is ``outside'' time and space, He
    lacks the temporal framework necessary for the act of ``deciding'' or
    ``causing,'' since causation presupposes temporal succession.
  \item \textbf{The problem of divine hiddenness.}
    If an omnibenevolent God exists, He would ensure that every sincere
    seeker could find adequate evidence for His existence.  The existence of
    reasonable non-belief is itself evidence against such a God
    (J.~L.~Schellenberg).
  \item \textbf{Mono\-theism's exclusive truth-claim.}
    Each monotheism declares all others false while offering the same
    \emph{type} of evidence (scripture, testimony, inner experience).  The
    symmetry of competing revelations refutes
    all equally.
\end{itemize}

%%  III — REVEALED RELIGIONS
\section{Revealed Religions: The Claim of Divine Revelation}

\subsection{Sheet 6: Epistemology of Revelation}

The concept of ``divine revelation'' faces several structural difficulties:

\begin{enumerate}
  \item \textbf{The authentication problem.}
    How does one distinguish a genuine divine revelation from a hallucination,
    a delusion, an epileptic seizure, or a
    political fabrication?  No non-circular criterion exists within the
    revelatory framework.
  \item \textbf{The transmission problem.}
    Even if an original revelation occurred, centuries of oral transmission,
    copying, editing, and translation introduce noise that make the current text an unreliable witness to the
    original event.
  \item \textbf{The interpretation problem.}
    Sacred texts are radically polysemic.  The same verses have been used to
    justify slavery and abolition, war and pacifism, monarchy
    and democracy.  If God intended a clear message, He failed.
  \item \textbf{The exclusivity problem.}
    Multiple religions claim mutually exclusive revelations.  Since each
    relies on the same \emph{kind} of evidence (a text said to be inspired),
    internal consistency alone cannot adjudicate between them.
\end{enumerate}

\subsection{Sheet 7: Miracles and Natural Law}

\begin{quote}
\textit{``No testimony is sufficient to establish a miracle, unless the
testimony be of such a kind, that its falsehood would be more miraculous,
than the fact, which it endeavours to establish.''}

\medskip
\hfill--- David Hume, \textit{Of Miracles} (1748)
\end{quote}

\begin{itemize}
  \item \textbf{Hume's maxim.}
    The probability of human error, exaggeration, or fraud \emph{always}
    exceeds the probability that a law of nature has been violated.
  \item \textbf{Survivor bias.}
    Miracle accounts survive because they are extraordinary.  The vast
    majority of prayers unanswered, of relics that healed nobody, vanish
    from the record.
  \item \textbf{Cultural dependence.}
    Miracle reports cluster in pre-scientific, low-literacy societies.
    They decline precisely where education, record-keeping, and medical
    knowledge advance.
\end{itemize}

%%  IV — JUDAISM
\section{Judaism: Historical and Theological Critique}

\subsection{Sheet 8: The Old Testament as History}

Modern biblical archaeology has
undermined the historicity of the foundational narratives:

\begin{itemize}
  \item \textbf{The Exodus.}
    No Egyptian record mentions a mass departure of slaves.
    Archaeological surveys of the Sinai find no trace of forty years of
    wandering by 600,000 men (plus families).  The story is best understood
    as a national origin myth crystallised during the Babylonian
    exile.
  \item \textbf{The Conquest of Canaan.}
    Jericho was unfortified and largely uninhabited at the
    supposed date of conquest.  The archaeological evidence suggests a gradual
    internal emergence of Israelite identity from Canaanite
    society, not a military invasion.
  \item \textbf{The United Monarchy.}
    The grandeur of Solomon's kingdom described in Kings is
    not supported by the modest material record of 10th-century~BCE
    Jerusalem.
\end{itemize}

\subsection{Sheet 9: Ethical Critique of the Hebrew Bible}

The moral content of the Torah presents acute difficulties for
any attempt to derive universal ethics from it:

\begin{itemize}
  \item \textbf{Commanded genocide (\emph{herem}).}
    Deuteronomy~20:16--17 and~1~Samuel~15:3 command the total extermination
    of conquered peoples, including women and children, explicitly attributed
    to God's will.
  \item \textbf{Slavery legislation.}
    Exodus~21 regulates slavery as a lawful institution.  Non-Israelite
    slaves may be owned permanently and bequeathed as property
    (Leviticus~25:44--46).
  \item \textbf{Collective punishment.}
    God punishes children for the sins of their parents
    (Exodus~20:5, Numbers~16:27--33).
  \item \textbf{Gender subordination.}
    Women are treated as property in marriage law (Deuteronomy~22:28--29),
    inheritance law, and purity regulations.
\end{itemize}

\paragraph{Hermeneutical Escape?}
The standard defence---``these laws must be read in historical
context''---concedes the central atheist point: the text is a human product
reflecting the moral limitations of its authors, not the utterance of a
timeless, perfect being.

\subsection{Sheet 10: The Chosen People Concept}

\begin{itemize}
  \item \textbf{Moral arbitrariness.}
    The election of one ethnic group as God's
    favoured people implies that the Creator of the universe has ethnic
    preferences---a claim that sits uneasily with any concept of universal
    justice.
  \item \textbf{From theology to politics.}
    The ``chosen people'' concept has been instrumentalised throughout
    history to justify territorial claims, cultural supremacism, and the
    exclusion of outsiders from moral concern.
  \item \textbf{Internal contradiction.}
    Later prophetic universalism (Isaiah~2:2--4,
    Micah~4:1--3) contradicts the particularism of Deuteronomic theology,
    exposing the composite, multi-authored nature of the tradition.
\end{itemize}

%%  V — CHRISTIANITY
\section{Christianity: Dogmas and Contradictions}

\subsection{Sheet 11: The Historical Jesus Problem}

\begin{itemize}
  \item \textbf{Sources.}
    The four canonical Gospels are not eyewitness
    accounts.  They were composed 40--70 years after the events they
    describe, in Greek, by authors drawing on earlier oral traditions and
    written sources (Q, Mark).
  \item \textbf{Contradictions.}
    The Gospels disagree on: the genealogy of Jesus (Matthew vs.\ Luke),
    the date of the crucifixion (synoptics vs.\ John), the events at the
    tomb, and the post-resurrection appearances.
  \item \textbf{Mythologisation.}
    Many elements of the Jesus narrative---virgin birth,
    star of Bethlehem, walking on water, resurrection---have parallels in
    pre-Christian Mediterranean religions (Dionysus,
    Osiris, Mithras), suggesting a pattern
    of mythological accretion rather than historical reporting.
  \item \textbf{Paul vs.\ Jesus.}
    Paul, the earliest Christian author, shows almost no interest in the
    teachings or biography of the earthly Jesus.  His ``Christ'' is a cosmic
    figure encountered in visions, not a remembered teacher.
\end{itemize}

\subsection{Sheet 12: The Trinity}

\begin{itemize}
  \item \textbf{Logical incoherence.}
    The doctrine asserts that God is three ``persons'' but one ``substance.''
    No definition of ``person'' and ``substance'' has ever been offered that
    makes this claim non-contradictory without evacuating both terms of
    meaning.
  \item \textbf{Historical contingency.}
    The Trinitarian formula was settled at Nicaea
    (325~CE) by majority vote under imperial pressure from
    Constantine.  Alternative Christologies
    (Arianism, Adoptionism,
    Modalism) were not refuted but \emph{suppressed}.
  \item \textbf{Biblical absence.}
    The word ``Trinity'' does not appear in the Bible.  The Johannine
    Comma (1~John~5:7--8), long cited as scriptural
    support, is a later interpolation absent from all early manuscripts.
\end{itemize}

\subsection{Sheet 13: Original Sin and Atonement}

\begin{itemize}
  \item \textbf{Inherited guilt.}
    The doctrine of Original Sin asserts that all humans
    are guilty because of Adam's transgression.  This violates the most
    elementary principle of justice: that guilt is personal and
    non-transferable.
  \item \textbf{The logic of atonement.}
    God requires a blood sacrifice---of Himself, to
    Himself---to forgive a transgression He could simply have pardoned.
    The mechanism is incoherent: how does the suffering of an innocent party
    erase the guilt of the guilty?
  \item \textbf{The ``sacrifice'' that costs nothing.}
    If Jesus is God and therefore cannot truly die, the ``supreme sacrifice''
    amounts to a weekend of discomfort followed by eternal glory---not a
    genuine sacrifice by any meaningful standard.
  \item \textbf{Moral hazard.}
    The doctrine of vicarious atonement teaches that personal responsibility
    can be transferred.  The deathbed conversion
    of a lifelong criminal is valued above the virtuous life of an
    unbeliever---a morally monstrous conclusion.
\end{itemize}

\subsection{Sheet 14: Eschatology: Hell and Eternal Punishment}

\begin{itemize}
  \item \textbf{Infinite punishment for finite sins.}
    No finite transgression can merit \emph{infinite} punishment.
    Eternal hell is not justice; it is sadism elevated to a cosmic principle.
  \item \textbf{The geography of injustice.}
    Billions of humans have lived and died without hearing the Christian
    message.  Condemning them for failing to believe what they never
    encountered is incompatible with any coherent concept of justice or
    mercy.
  \item \textbf{Fear as a tool of control.}
    The doctrine of hell functions not as a metaphysical truth but as a
    mechanism of psychological coercion, keeping believers obedient and
    sceptics silent.
\end{itemize}

%%  VI — THE PROTESTANT REFORMATION
\section{The Reformation: Grace, Faith, and Predestination}

\subsection{Luther and \textit{Sola Fide}}

\subsection{Sheet 15: Luther's Soteriological Revolution}

Martin Luther's break with Rome rested on a radical
reinterpretation of salvation:

\begin{enumerate}
  \item \textbf{\textit{Sola fide} (faith alone).}
    Humans are justified (declared righteous) not by works---moral effort,
    sacraments, pilgrimages---but solely by faith in Christ.
  \item \textbf{\textit{Sola gratia} (grace alone).}
    Even faith itself is not a human achievement but a gift of divine
    grace.  The human will is, in Luther's phrase, a
    ``\textit{servum arbitrium}''---a bound will,
    incapable of turning toward God without prior divine intervention.
  \item \textbf{\textit{Sola scriptura}.}
    The Bible alone is the supreme authority in matters of doctrine, against
    the Catholic triad of scripture, tradition, and papal magisterium.
\end{enumerate}

\paragraph{The Ethical Problem}
If salvation depends entirely on grace and not on moral conduct, then ethics
become instrumentally irrelevant to the ultimate human concern.  Luther
struggled with this implication and never resolved it satisfactorily.
His notorious advice to ``sin boldly'' (\textit{pecca
fortiter, sed fortius fide}) captures the paradox.

\subsection{Calvin and Double Predestination}

\subsection{Sheet 16: The Doctrine of Predestination}

John Calvin radicalised Luther's soteriology into the
doctrine of \textbf{double predestination}:

\begin{enumerate}
  \item \textbf{Unconditional election.}
    Before the creation of the world, God chose (\emph{elected}) some
    individuals for salvation and others for damnation.  This choice is
    entirely sovereign: it is not based on foreseen faith, merit, or any
    human quality.
  \item \textbf{Limited atonement.}
    Christ did not die for all humanity, but only for the elect.
    The ``universalist'' passages of the New Testament
    (John~3:16, 1~Timothy~2:4) must be reinterpreted to mean ``all kinds
    of people,'' not ``every individual.''
  \item \textbf{Irresistible grace.}
    Those whom God has chosen cannot refuse His grace.  Conversely, the
    reprobate cannot attain grace no matter how earnestly
    they seek it.
  \item \textbf{Perseverance of the saints.}
    The elect cannot fall from grace.  Apparent apostasy in a believer
    proves only that they were never truly elect.
\end{enumerate}

These five points are sometimes summarised by the acrostic
\textbf{TULIP}:\linebreak
\textbf{T}otal depravity,
\textbf{U}nconditional election,
\textbf{L}imited atonement,
\textbf{I}rresistible grace,
\textbf{P}erseverance of the saints.

\subsection{Sheet 17: Critique of Predestination}

\begin{itemize}
  \item \textbf{Moral monstrosity.}
    A God who creates beings \emph{for the express purpose} of damning them
    eternally is not merely unjust but indistinguishable from a cosmic
    torturer.  The traditional language of ``divine justice'' becomes an
    Orwellian inversion.

  \item \textbf{Destruction of moral agency.}
    If every human action is determined by prior divine decree, then praise
    and blame, reward and punishment, are meaningless.  The entire moral
    universe collapses into a puppet theatre.

  \item \textbf{Epistemic paralysis.}
    The doctrine generates intolerable anxiety: \emph{am I among the elect?}
    Since no human action can confirm or deny election, the believer is
    trapped in permanent uncertainty about the only question that
    (supposedly) matters.  Weber famously argued that this
    anxiety drove Calvinists toward compulsive worldly activity and
    capital accumulation as \emph{indirect} signs of election---the
    ``Protestant ethic'' as psychological
    compensation.

  \item \textbf{The problem of prayer.}
    If God's decrees are fixed from eternity, prayer is futile: it cannot
    change outcomes.  Yet Calvin insisted on its necessity---an incoherence
    masked by appeals to mystery.

  \item \textbf{Internal biblical contradiction.}
    Predestination sits in tension with numerous biblical passages that
    presuppose genuine human choice:
    ``Choose this day whom you will serve'' (Joshua~24:15);
    ``Whoever will, let him come'' (Revelation~22:17).
    Calvin's exegesis requires reading these as divine rhetoric addressed
    to beings who cannot, in fact, choose---a reading bordering on divine
    deception.
\end{itemize}

\subsection{Sheet 18: Grace: A Genealogy of the Concept}

\begin{itemize}
  \item \textbf{Pauline origins.}
    Paul's letters introduce \textit{charis} (grace) as
    God's unmerited favour, set against the ``works of the
    law''.  However, Paul's own usage is
    inconsistent: Romans~9 suggests absolute divine sovereignty, while
    Romans~2:6--7 implies judgement according to works.

  \item \textbf{Augustine's systematisation.}
    Augustine, shaped by his personal struggle with desire and his
    anti-Pelagian polemic, hardened grace into
    irresistible divine causality.  His doctrine of \textit{massa
    damnata}---the entire human race as a ``lump of
    damnation'' from which God rescues a few---laid the foundation for
    Calvin.

  \item \textbf{The Pelagian alternative.}
    Pelagius argued that humans have genuine free will
    and the natural capacity for moral perfection.  Grace is helpful but
    not indispensable.  His position was condemned---but its condemnation
    was a \emph{political} victory of the Augustinian party, not a
    philosophical refutation.

  \item \textbf{Modern assessment.}
    The concept of ``grace'' resolves nothing.  It is a placeholder for
    the mystery of why some believe and others do not---a mystery that
    naturalistic psychology (upbringing,
    temperament, social context) explains without supernatural residue.
\end{itemize}

\subsection{Synthesis: The Reformation as Radicalised Incoherence}

\subsection{Sheet 19: The Reformation's Deepened Contradictions}

Far from ``purifying'' Christianity, the Reformation \emph{intensified} its
deepest contradictions:

\begin{enumerate}
  \item \textbf{It replaced one authority with another.}
    \textit{Sola scriptura} did not liberate conscience; it substituted the
    authority of the pope with the authority of
    competing biblical interpreters.  The result was not freedom but
    fragmentation: tens of thousands of Protestant
    denominations, each claiming the
    ``plain sense'' of the same text.

  \item \textbf{It deepened the moral abyss.}
    Catholic theology, for all its faults, at least preserved a space for
    human moral effort.  Calvinist predestination annihilates it.  The
    Reformation's God is \emph{more} arbitrary, \emph{more} cruel, and
    \emph{more} indifferent to human striving than the medieval God it
    replaced.

  \item \textbf{It exposed the constructedness of doctrine.}
    The very fact that Christianity could be so radically reinterpreted---that
    ``the faith once delivered'' was, in practice, endlessly mutable---is
    itself powerful evidence that Christian doctrine is a human
    construction, not a
    fixed divine revelation.
\end{enumerate}

%%  CONCLUDING NOTE
\section*{Concluding Note}
\addcontentsline{toc}{section}{Concluding Note}

\begin{quote}
\textit{``We want to stand upon our own feet and look fair and square at the
world---its good facts, its bad facts, its beauties, and its ugliness; see
the world as it is and be not afraid of it.  Conquer the world by
intelligence and not merely by being slavishly subdued by the terror that
comes from it.''}

\medskip
\hfill--- Bertrand Russell, \textit{Why I Am Not a Christian} (1927)
\end{quote}

The purpose of these study sheets is not to promote hostility toward
believers, but to demonstrate that the arguments \emph{for} theism (and
especially for the specific doctrines of Judaism, Christianity, and
Protestantism) fail to meet the evidentiary and logical standards we
apply to every other domain of human inquiry.

Where evidence is absent, reason is the only honest guide.  Where
contradictions abound, intellectual courage demands that we name them.  And
where doctrines cause suffering---eternal damnation, predestined reprobation,
inherited guilt---moral decency demands that we reject them.

\endgroup
